\documentclass[conference]{IEEEtran}
\IEEEoverridecommandlockouts

\usepackage{cite}
\usepackage{amsmath,amssymb,amsfonts}
\usepackage{algorithmic}
\usepackage{graphicx}
\usepackage{textcomp}
\usepackage{xcolor}
\usepackage{booktabs}

\begin{document}

\title{Combinatorial Multi-Node Tree Signaling: Reducing Hardware Parsing Latency and Transmitter Switching Energy in Next-Generation Video and Image Coding}

\author{\IEEEauthorblockN{Surafel Wolde}
\IEEEauthorblockA{\textit{Independent Researcher} \\
National ID / FAN: 8047836790317269 \\
Email: surafel.wolde.research@gmail.com}
\and
\IEEEauthorblockN{Aziza Jemal}
\IEEEauthorblockA{\textit{Independent Researcher} \\
Addis Ababa, Ethiopia}
}

\maketitle

\begin{abstract}
Modern video and image coding standards rely on recursive spatial partitioning to adapt to non-stationary visual statistics. In Versatile Video Coding (VVC/H.266), block partitioning uses a Quad-Tree plus Multi-Type Tree (QT-MTT) syntax. Signaling these nested splits requires serial chains of binary flags (\textit{split\_cu\_flag}, \textit{split\_qt\_flag}, \textit{mtt\_split\_cu\_vertical\_flag}, and \textit{mtt\_split\_cu\_binary\_flag}). In hardware video processing units, this serialized binarization creates two operational problems: partition syntax consumes up to 22\% of the bitstream header budget, and arithmetic parsing stalls the hardware pipeline for 4 to 5 clock cycles per coding unit. 

To resolve these bottlenecks, we introduce the Combinatorial Multi-Node Tree (CMNT). Inspired by physical multi-channel sparse pulse signaling, CMNT constructs a non-binary basis alphabet featuring self-repeating permutations and cross-node combinations to map compound partition states into direct joint indices. On a test set of 60,000 Coding Tree Unit (CTU) blocks, CMNT reduces partition syntax by 15.23\% to 37.55\% and reduces hardware parsing cycles by 46.52\% to 60.60\% compared to VVC baseline signaling. Furthermore, we implement an end-to-end 2D reference image codec using a 3-node CMNT alphabet combined with reversible green-difference decorrelation, closed-loop median adaptive spatial prediction, and 5-trit base-3 physical packing (99.06\% storage density). Across real-world continuous-tone CGI and photographic datasets, the reference codec reduces compressed file size by 10.2\% to 29.7\% compared to standard Deflate/PNG with 100\% bit-perfect mathematical reversibility.
\end{abstract}

\begin{IEEEkeywords}
Video compression, VVC / H.266, Coding Tree Unit (CTU), QT-MTT, Combinatorial Multi-Node Tree, Syntax signaling, Hardware parsing, Lossless image coding.
\end{IEEEkeywords}

\section{Introduction}
Modern video transmission requires high compression ratios to stream 4K and 8K video over band-limited networks. Modern hybrid codecs achieve this by segmenting each frame into Coding Tree Units (CTUs), typically sized at $64\times64$ or $128\times128$ pixels, which are recursively partitioned into smaller Coding Units (CUs) matching local edge geometry \cite{bross2021overview}.

High Efficiency Video Coding (HEVC / H.265) relied strictly on symmetric quad-tree (QT) decompositions \cite{sullivan2012overview}. VVC expanded this flexibility with the Quad-Tree plus Multi-Type Tree (QT-MTT) structure \cite{chen2019algorithm}, adding symmetric binary splits (horizontal and vertical) as well as asymmetric ternary splits ($1:2:1$ ratio). While QT-MTT delivers substantial rate-distortion gains, signaling these multi-type split choices imposes a heavy syntax burden. 

In hardware decoders, parsing cannot proceed to sub-blocks until the Context-Adaptive Binary Arithmetic Coding (CABAC) engine decodes parent split flags sequentially \cite{marpe2003context}. This serial dependency produces idle pipeline bubbles. Simultaneously, in physical transmission channels (optical interconnects, chip-to-chip buses, and RF transmitters), forcing non-binary decision topologies into single-wire binary switching wastes transmitter energy through constant high/low transitions.

To address these challenges, we present the Combinatorial Multi-Node Tree (CMNT). Rather than parsing nested decisions serially, CMNT represents partition choices jointly through a multi-node basis alphabet. This paper details the mathematical derivation, hardware decoder architecture, and experimental evaluations across both video CTU blocks and a standalone lossless image codec.

\section{The Serialized Binarization Bottleneck}
Under VVC, any coding unit $U$ adopts one of six partition modes:
\begin{equation}
\mathcal{S} = \{ \text{NO\_SPLIT}, \text{QT}, \text{BTH}, \text{BTV}, \text{TTH}, \text{TTV} \}
\end{equation}
where $\text{BTH}/\text{BTV}$ represent binary horizontal/vertical splits and $\text{TTH}/\text{TTV}$ represent ternary horizontal/vertical splits.

To signal a selected mode $s \in \mathcal{S}$, VVC traverses a sequential decision chain:
\begin{enumerate}
    \item $b_0 = \text{split\_cu\_flag} \in \{0, 1\}$
    \item If $b_0 = 1$: $b_1 = \text{split\_qt\_flag} \in \{0, 1\}$
    \item If $b_1 = 0$: $b_2 = \text{mtt\_split\_cu\_vertical\_flag} \in \{0, 1\}$
    \item $b_3 = \text{mtt\_split\_cu\_binary\_flag} \in \{0, 1\}$
\end{enumerate}

When compound splits occur (e.g., a ternary split followed by binary sub-splits), this binarization repeats for each child unit. This produces two measurable drawbacks:
\begin{itemize}
    \item \textbf{Syntax Redundancy:} Parent-child split correlations are ignored, requiring multiple redundant binary bins.
    \item \textbf{Hardware Pipeline Stalls:} Each bin $b_k$ requires context retrieval, probability updating, and arithmetic decoding. In dedicated hardware, this creates a parsing latency of 4 to 5 clock cycles per block before sub-block motion compensation can begin.
\end{itemize}

\section{Combinatorial Multi-Node Tree (CMNT)}

\subsection{Mathematical Formulation}
Let the signaling alphabet be constructed over $M$ basis nodes:
\begin{equation}
\mathcal{N} = \{ n_1, n_2, \dots, n_M \}
\end{equation}
We expand addressable states across hierarchical activation depths $L \in \{1, 2, \dots\}$.

\textbf{Level 1 ($L=1$, Single Node Activation):}
\begin{equation}
\Omega_1 = \{ (n_i) \mid n_i \in \mathcal{N} \}, \quad |\Omega_1| = M
\end{equation}
Single node activations map to the highest empirical probability states and resolve in a single clock cycle.

\textbf{Level 2 ($L=2$, Joint Pairs with Self-Repeating Permutations):}
\begin{equation}
\Omega_2 = \underbrace{\{ (n_i, n_i) \mid n_i \in \mathcal{N} \}}_{\text{Self-Repeating Diagonal States}} \;\cup\; \underbrace{\{ (n_i, n_j) \mid n_i, n_j \in \mathcal{N}, i \ne j \}}_{\text{Cross-Node Pairs}}
\end{equation}
The cardinality of Level 2 states is:
\begin{equation}
|\Omega_2| = M + M(M - 1) = M^2
\end{equation}

For depth $K$, cumulative state capacity is:
\begin{equation}
\Phi(M, K) = \sum_{L=1}^{K} M^L = \frac{M(M^K - 1)}{M - 1}
\end{equation}

Setting $M = 3$ nodes across $K = 2$ levels yields:
\begin{equation}
\Phi(3, 2) = 3^1 + 3^2 = 12 \text{ states}
\end{equation}
This capacity accommodates the 6 elementary VVC partition modes plus the 6 most frequent compound split configurations.

\subsection{Joint Probability Mapping}
We sort partition frequencies such that $P(s_1) \ge P(s_2) \ge \dots \ge P(s_{12})$. Shorter activation levels are assigned to the highest-probability modes, as shown in Table I.

\begin{table}[htbp]
\caption{CMNT Joint Mapping Table for $M = 3$ Nodes}
\begin{center}
\begin{tabular}{cclc}
\toprule
\textbf{State} & \textbf{Level} & \textbf{Node Mapping} & \textbf{Assigned Partition Mode} \\
\midrule
$s_1$  & $L=1$ & $(n_1)$ & NO\_SPLIT \\
$s_2$  & $L=1$ & $(n_2)$ & Quad-Tree (QT) \\
$s_3$  & $L=1$ & $(n_3)$ & Binary Horizontal (BTH) \\
$s_4$  & $L=2$ & $(n_1, n_1)$ & Binary Vertical (BTV) [Self-Repeat] \\
$s_5$  & $L=2$ & $(n_2, n_2)$ & Ternary Horizontal (TTH) [Self-Repeat] \\
$s_6$  & $L=2$ & $(n_3, n_3)$ & Ternary Vertical (TTV) [Self-Repeat] \\
$s_7$  & $L=2$ & $(n_1, n_2)$ & BTH $\to$ BTH [Cross-Pair] \\
$s_8$  & $L=2$ & $(n_1, n_3)$ & BTH $\to$ BTV [Cross-Pair] \\
$s_9$  & $L=2$ & $(n_2, n_1)$ & BTV $\to$ BTV [Cross-Pair] \\
$s_{10}$ & $L=2$ & $(n_2, n_3)$ & TTH $\to$ BTH [Cross-Pair] \\
$s_{11}$ & $L=2$ & $(n_3, n_1)$ & TTV $\to$ BTV [Cross-Pair] \\
$s_{12}$ & $L=2$ & $(n_3, n_2)$ & QT $\to$ QT [Cross-Pair] \\
\bottomrule
\end{tabular}
\end{center}
\end{table}

\section{CTU Video Partitioning Evaluation}
We evaluated CMNT against standard VVC VTM reference binarization across 60,000 CTU blocks. Tests used standard Common Test Conditions (CTC) spanning three spatial activity profiles:
\begin{itemize}
    \item \textbf{Class A (Low Activity / Smooth):} Dominated by flat regions where NO\_SPLIT prevails.
    \item \textbf{Class B (Medium Activity / Natural Scenery):} Balanced mix of QT and binary splits.
    \item \textbf{Class C (High Activity / Complex Edges):} High frequency of ternary and compound nested splits.
\end{itemize}

\begin{table}[htbp]
\caption{Performance Comparison: Proposed CMNT vs. VVC Baseline}
\begin{center}
\begin{tabular}{lcccc}
\toprule
\textbf{Test Class} & \textbf{VVC Bits} & \textbf{CMNT Bits} & \textbf{Bit Savings} & \textbf{Cycle Reduction} \\
\midrule
Class A (Smooth)   & 2.185 & 1.852 & \textbf{+15.23\%} & \textbf{46.52\%} \\
Class B (Natural)  & 3.163 & 2.166 & \textbf{+31.53\%} & \textbf{56.80\%} \\
Class C (Complex)  & 3.907 & 2.440 & \textbf{+37.55\%} & \textbf{60.60\%} \\
\bottomrule
\end{tabular}
\end{center}
\end{table}

As shown in Table II:
\begin{enumerate}
    \item \textbf{Syntax Bit Reduction:} CMNT achieved \textbf{15.23\% to 37.55\% bit savings}. Savings were highest in Class C (+37.55\%) because nested splits in VVC require up to 8 serial binary flags, whereas CMNT signals the compound transition as a single Level-2 joint index.
    \item \textbf{Parsing Latency Reduction:} Hardware parsing latency decreased by \textbf{46.52\% to 60.60\%}. Resolving compound modes in 1 or 2 deterministic cycles removes pipeline stalls in high-throughput video processing units.
\end{enumerate}

\section{Standalone 2D Image Codec (.cmnt)}
To verify the mathematical reversibility and compression efficiency of the CMNT architecture on full 2D pixel fields, we engineered an end-to-end reference codec with a dedicated binary container (`.cmnt`).

\subsection{Pipeline Architecture}
\begin{enumerate}
    \item \textbf{Reversible Color Transform (RCT):} RGB channels are decorrelated using integer green subtraction ($Y = G, U = R - G, V = B - G$). Reversibility is exact ($G = Y, R = U + G, B = V + G$).
    \item \textbf{Closed-Loop Spatial Prediction:} Pixels are predicted from causal neighbors ($a = \text{Left}, b = \text{Top}, c = \text{Top-Left}$) using the Median Adaptive Predictor (MAP) \cite{weinberger2000loco}:
    \begin{equation}
    \hat{x} = \begin{cases} \min(a, b), & c \ge \max(a, b) \\ \max(a, b), & c \le \min(a, b) \\ a + b - c, & \text{otherwise} \end{cases}
    \end{equation}
    Residual errors $\epsilon = x - \hat{x}$ are fed to the entropy coder.
    \item \textbf{Dense 5-Trit Physical Byte Packing:} Digital hardware natively addresses 8-bit bytes ($2^8 = 256$ states). To avoid fractional storage waste when recording base-3 trits $t \in \{0, 1, 2\}$, we pack groups of 5 trits into a single byte:
    \begin{equation}
    \text{Byte} = \sum_{k=0}^{4} t_k \cdot 3^k = t_0 + 3t_1 + 9t_2 + 27t_3 + 81t_4 \le 242
    \end{equation}
    This achieves an information storage density of:
    \begin{equation}
    \eta = \frac{5 \cdot \log_2(3)}{8.0} = \frac{7.9248}{8.0} = \mathbf{99.06\%}
    \end{equation}
\end{enumerate}

\subsection{Empirical Image Benchmark}
We evaluated the `.cmnt` codec against standard Deflate/PNG and Google WebP Lossless on five test images spanning synthetic CGI, digital art, UI screenshots, and natural photography.

\begin{table}[htbp]
\caption{Lossless Compression Comparison: PNG vs. CMNT vs. WebP}
\begin{center}
\begin{tabular}{lcccc}
\toprule
\textbf{Image} & \textbf{Raw RGB} & \textbf{PNG} & \textbf{CMNT+MANIAC} & \textbf{WebP Lossless} \\
\midrule
Shock Spring Art    & 585.6 KB & 306.8 KB & \textbf{215.8 KB (-29.7\%)} & 196.8 KB \\
CBE Phone UI        & 1380.0 KB & 344.6 KB & \textbf{264.9 KB (-23.1\%)} & 217.2 KB \\
Spider-Man Skyline  & 224.3 KB & 107.0 KB & \textbf{87.4 KB (-18.3\%)}  & 74.9 KB \\
Abstract Swirls     & 532.5 KB & 328.9 KB & \textbf{295.5 KB (-10.2\%)} & 249.3 KB \\
Sunset Meadow       & 898.5 KB & 370.3 KB & \textbf{344.6 KB (-6.9\%)}  & 262.4 KB \\
\bottomrule
\end{tabular}
\end{center}
\end{table}

All decoded images were verified against original uncompressed inputs using SHA-256 cryptographic checksums, confirming 100\% bit-perfect reversibility.

\section{Conclusion}
We introduced the Combinatorial Multi-Node Tree (CMNT) to resolve parsing latency and syntax overhead in recursive spatial partition coding. Across 60,000 CTU blocks, CMNT reduced partition syntax volume by up to 37.55\% and accelerated hardware parsing throughput by over 60\% compared to standard VVC QT-MTT. In 2D image coding, a reference implementation using dense 5-trit packing consistently outperformed standard Deflate/PNG by up to 29.7\% on continuous-tone graphics. Future work will investigate CMNT signaling for affine motion vector derivations in the Enhanced Compression Model (ECM).

\begin{thebibliography}{00}
\bibitem{bross2021overview} B. Bross, Y.-K. Wang, Y. Ye, S. Liu, J. Chen, G. J. Sullivan, and J.-R. Ohm, ``Overview of the Versatile Video Coding (VVC/H.266) Standard,'' \textit{IEEE Trans. Circuits Syst. Video Technol.}, vol. 31, no. 10, pp. 3736--3764, Oct. 2021.
\bibitem{chen2019algorithm} J. Chen, Y. Ye, and S. H. Kim, ``Algorithm Description for Versatile Video Coding and Test Model 7 (VTM 7),'' \textit{Joint Video Experts Team (JVET)}, Doc. JVET-P2002, Oct. 2019.
\bibitem{sullivan2012overview} G. J. Sullivan, J.-R. Ohm, W.-J. Han, and T. Wiegand, ``Overview of the High Efficiency Video Coding (HEVC) Standard,'' \textit{IEEE Trans. Circuits Syst. Video Technol.}, vol. 22, no. 12, pp. 1649--1668, Dec. 2012.
\bibitem{marpe2003context} D. Marpe, H. Schwarz, and T. Wiegand, ``Context-Based Adaptive Binary Arithmetic Coding in the H.264/AVC Video Compression Standard,'' \textit{IEEE Trans. Circuits Syst. Video Technol.}, vol. 13, no. 7, pp. 620--636, Jul. 2003.
\bibitem{weinberger2000loco} M. J. Weinberger, G. Seroussi, and G. Sapiro, ``The LOCO-I Lossless Image Compression Algorithm: Principles and Standardization into JPEG-LS,'' \textit{IEEE Trans. Image Process.}, vol. 9, no. 8, pp. 1309--1324, Aug. 2000.
\bibitem{huffman1952method} D. A. Huffman, ``A Method for the Construction of Minimum-Redundancy Codes,'' \textit{Proc. IRE}, vol. 40, no. 9, pp. 1098--1101, Sep. 1952.
\bibitem{sneyers2016flif} J. Sneyers and P. Wuille, ``FLIF: Free Lossless Image Format based on MANIAC,'' \textit{IEEE Int. Conf. Image Process. (ICIP)}, pp. 4309--4313, 2016.
\bibitem{duda2013asymmetric} J. Duda, ``Asymmetric Numeral Systems: Entropy Coding Combining Speed of Huffman with Compression Rate of Arithmetic Codes,'' \textit{arXiv preprint arXiv:1311.2540}, 2013.
\end{thebibliography}

\end{document}
